\subsection{DEFINITION OF A FILTER}

DEFINITION 1. A filter on a set X is a set F of subsets of X which has the following properties:

(F \(_{I}\) ) Every subset of X which contains a set of F belongs to F.

\((\mathbf{F}_{\mathrm{II}})\) Every finite intersection of sets of \(\mathfrak{F}\) belongs to \(\mathfrak{F}\) .

\((\mathbf{F}_{\mathrm{III}})\) The empty set is not in \(\mathfrak{F}\) .

It follows from \((\mathbf{F}_{\mathbf{II}})\) and \((\mathbf{F}_{\mathbf{III}})\) that every finite intersection of sets of F is non-empty.

A filter \(\mathfrak{F}\) on X defines a structure on X, the axioms of which are \((\mathrm{F_I})\) , \((\mathrm{F_{II}})\) and \((\mathrm{F_{III}})\) ; this structure is called a structure of a filtered set, and the set X endowed with this structure is called a set filtered by \(\mathfrak{F}\) .

Axiom \((\mathbf{F}_{\mathrm{II}})\) is equivalent to the conjunction of the following two axioms: \((\mathbf{F}_{\mathrm{II}a})\) The intersection of two sets of F belongs to F.

\((\mathbf{F}_{\mathrm{II}b})\) X belongs to \(\mathfrak{F}\) .

Axioms \((\mathbf{F}_{\mathbf{II}\ \mathbf{b}})\) and \((\mathbf{F}_{\mathbf{III}})\) show that there is no filter on the empty set.

In order for a set of subsets which satisfies \((\mathbf{F}_{\mathbf{I}})\) also to satisfy \((\mathbf{F}_{\mathbf{II}\mathbf{b}})\) it is necessary and sufficient that it is not empty. A set of subsets which satisfies \((\mathbf{F}_{\mathbf{I}})\) also satisfies \((\mathbf{F}_{\mathbf{III}})\) if and only if it is different from \(\mathfrak{P}(\mathbf{X})\) .

Examples of filters. 1) If \(\mathbf{X} \neq \emptyset\) , the set of subsets consisting of \(\mathbf{X}\) alone is a filter on \(\mathbf{X}\) . More generally, the set of all subsets of \(\mathbf{X}\) which contain a given non-empty subset \(\mathbf{A}\) of \(\mathbf{X}\) is a filter on \(\mathbf{X}\) .

\begin{enumerate}
\def\labelenumi{\arabic{enumi})}
\setcounter{enumi}{1}
\item
  In a topological space X, the set of all neighbourhoods of an arbitrary non-empty subset A of X (and in particular the set of all neighbourhoods of a point of X) is a filter, called the neighbourhood filter of A.
\item
  If X is an infinite set, the complements of the finite subsets of X are the elements of a filter. The filter of complements of finite subsets of the set N of integers \(\geqslant0\) is called the Fréchet filter.
\end{enumerate}

